% ===========================================================================
% BACKGROUND AND RELATED WORK --- argument skeleton (skill S14).
% Organize by TECHNICAL DISTINCTION, not one-paragraph-per-paper. For each
% category state: protected assets, public/private weight assumption, TCB,
% nonlinear handling, autoregressive support, training support, formal security
% level, difference from this work. Do NOT use related work to make unsupported
% novelty claims. ALL citations are \citeneeded until verified (see
% audit/citation_inventory.md). Prior draft flagged its bib as placeholders.
% ===========================================================================
\section{Background and Related Work}
\label{sec:related}

% ---- CENTRAL QUESTION ------------------------------------------------------
% Where does obfuscated-execution-for-public/private-weight generative LLMs sit
% relative to TEE-offload, matrix obfuscation, crypto inference, and attacks?

% ---- 3.0 Background (brief) ------------------------------------------------
% Decoder-only Transformer, RoPE, GQA/MQA, RMSNorm/LayerNorm, SwiGLU, KV-cache,
% LoRA. \citeneeded{Transformer}\citeneeded{RoPE}\citeneeded{GQA}\citeneeded{RMSNorm}\citeneeded{LoRA}

% ---- 3.1 TEE-assisted NN inference / hybrid TEE+accelerator ----------------
% Slalom (linear offload w/ verification), Graviton, TensorSCONE, Chiron.
% Difference: those target CNNs / non-autoregressive; we target decoder LLM
% generation + LoRA training with a masked accelerator path.
% \citeneeded{Slalom}\citeneeded{Graviton}\citeneeded{TEE inference}
% \TODO{Verify each venue/year from original papers (citation_inventory CI-*).}

% ---- 3.2 Algebraic / secret linear-transform obfuscation -------------------
% STIP, ObfuscaTune, Amulet-style right-mask nonlinear islands, permutation/
% Arrow matrix obfuscation. TABLE: protected assets / public-vs-private weights
% / TCB / nonlinear handling / autoregressive / training / formal level.
% \citeneeded{STIP}\citeneeded{ObfuscaTune}\citeneeded{Amulet}\citeneeded{Arrow permutation}
% Difference: \TODO{state crisply what is new vs Amulet -- NOV-01 reviewer risk;
%   KV-append breaks left masks, dense masks don't commute with nonlinears,
%   LoRA rank/gradient leakage. Must not read as "Amulet + LoRA".}

% ---- 3.3 Cryptographic private inference (FHE / MPC / hybrid) ---------------
% CryptoNets, GAZELLE, Delphi, SecureML, MiniONN, BumbleBee, CONJFORMER.
% Position as COMPLEMENTARY (exact, stronger guarantees, higher nonlinear cost),
% NOT replaced by our work. \citeneeded{BumbleBee}\citeneeded{CONJFORMER}\citeneeded{Delphi}

% ---- 3.4 KV-cache privacy --------------------------------------------------
% \citeneeded{KV-cache privacy} \TODO{find concrete references.}

% ---- 3.5 Private / protected parameter-efficient fine-tuning ---------------
% Protected LoRA training, split-learning gradient protection, DP-SGD.
% \citeneeded{private LoRA training}\citeneeded{DP-SGD}

% ---- 3.6 Attack & leakage analysis (motivation for our threat model) -------
% Activation/embedding inversion (EIA), gradient leakage
% (DLG/iDLG/GradInversion), membership inference, BRE, PIA. Present as the
% attacker toolbox our leakage analysis is evaluated against.
% \citeneeded{EIA}\citeneeded{DLG}\citeneeded{BRE}\citeneeded{PIA}\citeneeded{membership inference}

% ---- REVIEWER RISKS --------------------------------------------------------
% \risk{NOV-01/PWR-01: novelty vs Amulet + fair "different problem" framing.}
% \risk{Entire prior-draft bibliography was placeholder -- verify every entry.}
\TODO{Write related-work prose after citation\_inventory.md verification pass; use the comparison table as the backbone.}
